\documentclass[es]{../../elegantbook}
% ========== Portada e información básica ==========
\title{Documento de prueba ElegantBook}
\subtitle{Características principales de la plantilla}
\author{Ingeniero de prueba}
\institute{Grupo de prueba de plantillas}
\date{\today}
\version{v1.0}
\extrainfo{Este documento solo sirve para probar funciones de plantilla y no tiene valor académico}
\bioinfo{Lote de prueba}{TEST-FULL-20260506}
\logo{../../figure/logo-blue.png}
\cover{../../figure/cover.jpg}

% ========== Índice y configuraciones básicas ==========
\setcounter{tocdepth}{3} % Mostrar índice hasta nivel 3

% ========== Referencias bibliográficas de prueba ==========
\addbibresource[location=local]{../en/reference.bib}

\begin{document}

\maketitle
\frontmatter
\tableofcontents

% ========== Prólogo (con secciones) ==========
\pagenumbering{Roman}
\setcounter{page}{0}
\chapter*{Prólogo de prueba}
\markboth{Prólogo de prueba}{}
\section*{Descripción del documento}
\markright{Descripción del documento}
Este documento es un caso de prueba completo para ElegantBook, cubriendo todas las funciones básicas y avanzadas de la plantilla. Todo el contenido son datos de prueba ficticios sin significado práctico.
\newpage
\section*{Alcance de la prueba}
\markright{Alcance de la prueba}
Esta prueba incluye: portada, índice, fuentes, listas multinivel, fórmulas complejas, todos los entornos matemáticos, referencias cruzadas, figuras y tablas, notas marginales, resúmenes de capítulo, ejercicios, bibliografía y apéndices.

\mainmatter

% ========== Capítulo 1: Maquetación básica y listas multinivel ==========
\chapter{Prueba de maquetación y listas}
\begin{introduction}[Ítems de prueba de este capítulo]
\item Fuentes chinas y formato de texto
\item Anidamiento de lista desordenada de 3 niveles
\item Anidamiento de lista ordenada de 3 niveles
\item Párrafos normales y texto enfatizado
\end{introduction}

\section{Prueba de fuente y texto}
Prueba de texto normal. \textbf{Negrita}, \textit{Cursiva}, \underline{Subrayado}, \textcolor{red}{Prueba de texto rojo}.

\section{Prueba de lista desordenada de 3 niveles}
\begin{itemize}
\item Ítem de prueba nivel 1 A
\item Ítem de prueba nivel 1 B
    \begin{itemize}
    \item Ítem de prueba nivel 2 B1
    \item Ítem de prueba nivel 2 B2
        \begin{itemize}
        \item Ítem de prueba nivel 3 B2-1
        \item Ítem de prueba nivel 3 B2-2
            \begin{itemize}
            \item Ítem de prueba nivel 4 B2-1a
            \item Ítem de prueba nivel 4 B2-2b
            \end{itemize}
        \end{itemize}
    \end{itemize}
\end{itemize}

\section{Prueba de lista ordenada de 3 niveles}
\begin{enumerate}
\item Prueba ordenada nivel 1 1
\item Prueba ordenada nivel 1 2
    \begin{enumerate}
    \item Prueba ordenada nivel 2 2-1
    \item Prueba ordenada nivel 2 2-2
        \begin{enumerate}
        \item Prueba ordenada nivel 3 2-2-1
        \item Prueba ordenada nivel 3 2-2-2
            \begin{enumerate}
            \item Prueba ordenada nivel 4 2-2-2-1
            \item Prueba ordenada nivel 4 2-2-2-2
            \end{enumerate}
        \end{enumerate}
    \end{enumerate}
\end{enumerate}

% ========== Capítulo 2: Prueba de fórmulas matemáticas complejas ==========
\chapter{Prueba de fórmulas matemáticas complejas}
\begin{introduction}
\item Fórmula compleja de una sola línea
\item Fórmulas multilínea (align)
\item Matrices / Determinantes
\item Integrales múltiples / Sumas / Límites / Fracciones
\item Etiquetas de fórmula y referencias cruzadas
\end{introduction}

\section{Fórmula compleja de una sola línea}
Fórmulas complejas en línea: $\sum_{i=1}^n \frac{x_i^2}{\sqrt{y_i+1}} \geq \bar{x}^2$, $\lim_{n\to\infty} \left(1+\frac{1}{n}\right)^n = e$.

\section{Fórmulas multilínea y matrices}
\begin{align}
a^2 + b^2 &= c^2 \label{eq:gougu} \\
\int_{0}^{1}\int_{0}^{1} f(x,y) dxdy &= 1 \label{eq:double} \\
\begin{vmatrix}
a & b \\
c & d
\end{vmatrix} &= ad-bc \label{eq:det}
\end{align}

\section{Fórmulas con operadores grandes}
\begin{equation}\label{eq:series}
\sum_{k=0}^{\infty} \frac{x^k}{k!} = e^x,\quad \oint_{C} Pdx+Qdy = \iint_{D}\left(\frac{\partial Q}{\partial x}-\frac{\partial P}{\partial y}\right)dxdy
\end{equation}

Referencia de fórmula: Teorema de Pitágoras \eqref{eq:gougu}, integral doble \eqref{eq:double}, determinante \eqref{eq:det}, serie \eqref{eq:series}.

% ========== Capítulo 3: Prueba completa de entornos matemáticos ==========
\chapter{Prueba completa de entornos matemáticos}
\begin{introduction}
\item Definición / Teorema / Lema / Corolario / Axioma / Postulado
\item Proposición / Ejemplo / Problema / Ejercicio
\item Nota / Observación / Conclusión / Suposición / Propiedad
\item Entorno de solución / demostración
\end{introduction}

\section{Prueba de entornos}

% 1. Definición
\ifdefined\simpleTest
\begin{definition}[Definición de prueba]\label{def:test}
\else
\begin{definition}[Definición de prueba]{test}
\fi
Este es el contenido de prueba del entorno definición, sin significado matemático real, solo para verificación de estilo.
\end{definition}

% 2. Teorema
\ifdefined\simpleTest
\begin{theorem}[Teorema de prueba]\label{thm:test}
\else
\begin{theorem}[Teorema de prueba]{test}
\fi
Si $x>0$, entonces $x+1>1$, correspondiente a la fórmula \eqref{eq:gougu}.
\end{theorem}

% 3. Lema
\ifdefined\simpleTest
\begin{lemma}[Lema de prueba]\label{lem:test}
\else
\begin{lemma}[Lema de prueba]{test}
\fi
Para cualquier número real $x$, se cumple $x^2\geq0$.
\end{lemma}

% 4. Corolario
\ifdefined\simpleTest
\begin{corollary}[Corolario de prueba]\label{cor:test}
\else
\begin{corollary}[Corolario de prueba]{test}
\fi
Del Lema \ref{lem:test} se deduce: $(x-1)^2\geq0$.
\end{corollary}

% 5. Axioma
\ifdefined\simpleTest
\begin{axiom}[Axioma de prueba]\label{axi:test}
\else
\begin{axiom}[Axioma de prueba]{test}
\fi
Axioma de prueba, sin significado práctico.
\end{axiom}

% 6. Postulado
\ifdefined\simpleTest
\begin{postulate}[Postulado de prueba]\label{pos:test}
\else
\begin{postulate}[Postulado de prueba]{test}
\fi
Postulado de prueba, solo para verificación de entorno.
\end{postulate}

% 7. Proposición
\ifdefined\simpleTest
\begin{proposition}[Proposición de prueba]\label{pro:test}
\else
\begin{proposition}[Proposición de prueba]{test}
\fi
Contenido de proposición de prueba, asociado a la Definición \ref{def:test}.
\end{proposition}

% 8. Ejemplo
\begin{example}\label{exa:test}
Prueba de entorno ejemplo, numeración automática.
\end{example}

% 9. Problema
\begin{problem}\label{probl:test}
Prueba de entorno problema, para plantear cuestiones por resolver.
\end{problem}

% 10. Ejercicio
\begin{exercise}\label{exer:test}
Calcular: $1+2+3+\dots+10$.
\end{exercise}

% 11. Nota
\begin{note}
Entorno nota: sin numeración, para recordatorios clave.
\end{note}

% 12. Observación
\begin{remark}
Entorno observación: para explicaciones complementarias.
\end{remark}

% 13. Conclusión
\begin{conclusion}
Entorno conclusión: contenido resumido de prueba.
\end{conclusion}

% 14. Suposición
\begin{assumption}
Entorno suposición: configuración de condiciones de prueba.
\end{assumption}

% 15. Propiedad
\begin{property}
Entorno propiedad: características del objeto de prueba.
\end{property}

% 16. Solución
\begin{solution}
Solución al Ejercicio \ref{exer:test}: La suma vale $55$ (respuesta de prueba).
\end{solution}

% 17. Demostración
\begin{proof}
Demostración del Teorema \ref{thm:test}: Como $x>0$, sumando 1 a ambos lados queda demostrado.
\end{proof}

% ========== Capítulo 4: Referencias cruzadas · Figuras · Notas marginales ==========
\chapter{Referencias cruzadas y funciones extendidas}
\begin{introduction}
\item Todos los tipos de referencias cruzadas
\item Prueba sencilla de figuras y tablas
\item Prueba de función de nota marginal
\item Verificación de navegación entre capítulos
\end{introduction}

\section{Resumen de todas las referencias cruzadas}
Definición \ref{def:test}, Teorema \ref{thm:test}, Lema \ref{lem:test}, Corolario \ref{cor:test}, Axioma \ref{axi:test}, Postulado \ref{pos:test}, Proposición \ref{pro:test}, Ejemplo \ref{exa:test}, Problema \ref{probl:test}, Ejercicio \ref{exer:test}.

\section{Prueba de figuras y tablas}
\begin{figure}[h]
    \centering
    \includegraphics[width=0.5\textwidth]{../../image/scatter.jpg}
    \caption{Figura de prueba}\label{fig:test}
\end{figure}
\begin{table}[h]
    \centering
    \begin{tabular}{c*{6}{c}}
    \toprule
        $n$ & 1 & 2 & 3 & 4 & 5 & 6 \\
    \midrule
        $a_n$ & 1 & 1 & 2 & 3 & 5 & 8 \\ 
    \bottomrule
    \end{tabular}
    \caption{Tabla de prueba}\label{tab:test}
\end{table}
Referencia de figura/tabla: \figref{fig:test} es la figura de prueba, \tabref{tab:test} es la tabla de prueba.

% ========== Ejercicios al final del capítulo ==========
\begin{problemset}[Ejercicios de prueba del capítulo]
\item Verificar el estilo de visualización de la Definición \ref{def:test}
\item Deduzca la fórmula \eqref{eq:double}
\item Comprobar si la Figura \ref{fig:test} se muestra correctamente
\end{problemset}

% ========== Bibliografía ==========
\nocite{*}
\printbibliography

% ========== Apéndice ==========
\appendix
\chapter{Apéndice de prueba}
El contenido del apéndice sirve para verificar capítulos de apéndice, encabezados, numeración de páginas y estilo de maquetación, sin información real.
\section{Apéndice de prueba 1}
El contenido del apéndice sirve para verificar capítulos de apéndice, encabezados, numeración de páginas y estilo de maquetación, sin información real.
\newpage
\section{Apéndice de prueba 2}
El contenido del apéndice sirve para verificar capítulos de apéndice, encabezados, numeración de páginas y estilo de maquetación, sin información real.
\end{document}